% Pregunta 574 (2016-noviembre-OEP2016, n.º 19): restricción de liquidez y equivalencia ricardiana.
% Recta de balance intertemporal de pendiente -(1+r) con la dotación E. Sin poder endeudarse, el tramo
% a la derecha de E es inaccesible y el consumidor impaciente se queda en E en vez de en A.
% Una bajada de impuestos hoy financiada con deuda deja la recta igual pero lleva la dotación a E':
% el consumidor restringido consume más hoy (la deuda no es neutral); uno sin restricción seguiría en A.
% Autor: Víctor Gutiérrez Marcos
\documentclass[tikz,border=6pt]{standalone}
\input{tikz-tcee}
\begin{document}
\begin{tikzpicture}[tcee]
  % Recta: C2 = 4.8 - 1.1 C1 (1+r = 1,1)
  \def\pend{1.1}
  % Conjunto alcanzable con la restricción inicial (C1 <= 1,4)
  \fill[area] (0,0) -- (0,4.8) -- (1.4,3.26) -- (1.4,0) -- cycle;
  % Lo que se gana con la rebaja (C1 entre 1,4 y 2,4)
  \fill[area2] (1.4,0) -- (1.4,3.26) -- (2.4,2.16) -- (2.4,0) -- cycle;
  \ejes{$C_1$}{$C_2$}{6.2}{5.4}
  % Recta de balance: accesible hasta E, inaccesible después
  \draw[curva1] (0,4.8) -- (1.4,3.26);
  \draw[curva1,dashed,opacity=0.55] (1.4,3.26) -- (4.364,0);
  \node[etiqueta,text=tceemorado,anchor=south west] at (0.8,4.2) {pendiente $-(1+r)$};
  % Fronteras verticales: antes en Y1-T1, después en Y1-T1'
  \draw[curva1] (1.4,3.26) -- (1.4,0);
  \draw[curvanueva,desplazada] (2.4,2.16) -- (2.4,0);
  % Curvas de indiferencia del consumidor impaciente
  \draw[curva3,dashed,domain=2.15:5.6,samples=60] plot (\x,{1.28*(3.2/\x)^2.75});
  \draw[curva3,thin,domain=1.33:4.4,samples=60] plot (\x,{3.26*(1.4/\x)^2.75});
  \draw[curva3,domain=1.95:5.4,samples=60] plot (\x,{2.16*(2.4/\x)^2.75});
  % Desplazamiento de la dotación
  \draw[flecha] (1.55,3.30) -- (2.35,2.42);
  % Puntos
  \node[punto,label={[etiqueta]above right:$E$}] at (1.4,3.26) {};
  \node[punto,fill=tceerojo,label={[etiqueta,text=tceerojo]below left:$E'$}] at (2.4,2.16) {};
  \node[punto,fill=tceeverde,label={[etiqueta,text=tceeverde]below left:$A$}] at (3.2,1.28) {};
  % Proyecciones
  \draw[guia] (1.4,3.26) -- (0,3.26) node[left] {$Y_2-T_2$};
  \draw[guia] (2.4,2.16) -- (0,2.16) node[left,text=tceerojo] {$Y_2-T_2'$};
  \node[etiqueta,below] at (1.4,0) {$Y_1-T_1$};
  \node[etiqueta,below,text=tceerojo] at (2.4,-0.42) {$Y_1-T_1'$};
  % Notas
  \node[etiqueta,font=\footnotesize,align=left,text=tceeverde,anchor=west] at (3.5,3.6) {$A$: óptimo sin restricción\\(no cambia con la rebaja)};
  \node[etiqueta,font=\footnotesize,align=left,anchor=west,text=tceemorado] (nota) at (4.15,1.85) {tramo inaccesible:\\exigiría endeudarse};
  \draw[->,thin,tceemorado] (4.75,1.45) -- (4.2,0.26);
\end{tikzpicture}
\end{document}
